
\documentclass{article}
\usepackage{colortbl}
\usepackage{makecell}
\usepackage{multirow}
\usepackage{supertabular}

\begin{document}

\newcounter{utterance}

\centering \large Interaction Transcript for game `eqbench', experiment `eq\_bench', episode 17 with reflexion{-}Qwen3.5{-}9B.
\vspace{24pt}

{ \footnotesize  \setcounter{utterance}{1}
\setlength{\tabcolsep}{0pt}
\begin{supertabular}{c@{$\;$}|p{.15\linewidth}@{}p{.15\linewidth}p{.15\linewidth}p{.15\linewidth}p{.15\linewidth}p{.15\linewidth}}
    \# & \multicolumn{2}{c}{Player} && \multicolumn{2}{c}{Game Master} \\
    \hline

    \theutterance \stepcounter{utterance}  
    & & & \multicolumn{4}{p{0.6\linewidth}}{
        \cellcolor[rgb]{0.9,0.9,0.9}{
            \makecell[{{p{\linewidth}}}]{
                \texttt{\tiny{[P1$\langle$GM]}}
                \texttt{You have attempted this same task before and did not succeed. These are the plans you wrote after those attempts.} \\
\texttt{Trial \#1: New plan: I failed because I assigned a high score to "Triumphant" (6) despite Grayson's aggressive, paranoid, and ultimately isolated rhetoric, which suggests he feels defensive rather than victorious. His accusations are met with cold scientific rebuttals, not validation, leaving him feeling unheard and cornered. I should have recognized that his "evidence" was dismissed as "stories," leading to a sense of futility and defensiveness rather than triumph. Next time, I will lower the "Triumphant" score to 0 or 1, increase "Worried" to reflect his fear of being exposed as a liar, and assign the highest score (9 or 10) to "Justified" to capture his deep{-}seated belief that he is the only one seeing the truth, even if he is socially isolated.} \\
\\ 
\texttt{Trial \#2: New plan: I failed because I assigned a high score to "Triumphant" (6) and "Worried" (6), misinterpreting Grayson's aggressive rhetoric as a mix of victory and anxiety. In reality, Grayson's dialogue is a monologue of unyielding conviction; he dismisses all counter{-}arguments as "stories" and "paranoia," indicating he feels morally and intellectually superior to Sinclair. He does not feel triumphant because he hasn't won the argument, nor worried because he believes he holds the ultimate truth. The correct strategy is to recognize his deep{-}seated sense of moral righteousness despite the lack of external validation. I will now assign the highest score (10) to "Justified" to reflect his unwavering belief in his cause, and set "Triumphant" to 0, as his tone is one of righteous indignation, not victory. "Inferior" remains 0, and "Worried" drops to 1, reflecting only a slight concern about being labeled a liar, which is secondary to his primary feeling of being right.} \\
\\ 
\texttt{Justified: 10} \\
\texttt{Inferior: 0} \\
\texttt{Triumphant: 0} \\
\texttt{Worried: 1} \\
\\ 
\\ 
\\ 
\texttt{Your task is to predict the likely emotional responses of a character in this dialogue:} \\
\\ 
\texttt{Grayson: You know, Doc, I've been watching your work for years. All that time and money spent on Mars, and what have you got to show for it? } \\
\texttt{Dr. Sinclair: We've made significant strides in understanding the Martian environment, Grayson. It's not about instant results. } \\
\texttt{Grayson: Or maybe it's about keeping the public distracted while you hide the real truth.} \\
\texttt{Dr. Sinclair: And what truth would that be, Mr. Hawk?} \\
\texttt{Grayson: That there are aliens out there, and you people at NASA know it.} \\
\texttt{Dr. Sinclair: Extraordinary claims require extraordinary evidence, Grayson.} \\
\texttt{Grayson: Oh, I have evidence. I've got testimonials, photos, documents...} \\
\texttt{Dr. Sinclair: Anecdotes and blurry photos aren't evidence, Grayson. They're stories. } \\
\texttt{Grayson: Stories that you're too afraid to believe.} \\
\texttt{Dr. Sinclair: No, stories that don't hold up under scientific scrutiny. } \\
\texttt{Grayson: You're just a puppet, Sinclair. They've got you so convinced of your 'facts' that you can't see the truth.} \\
\texttt{Dr. Sinclair: And you're a sensationalist, Hawk. You've built a career on paranoia and misinformation.} \\
\texttt{Grayson: At least I'm not lying to the public.} \\
\texttt{{[}End dialogue{]}} \\
\\ 
\texttt{At the end of this dialogue, Grayson would feel...} \\
\texttt{Justified} \\
\texttt{Inferior} \\
\texttt{Triumphant} \\
\texttt{Worried} \\
\\ 
\texttt{Give each of these possible emotions a score from 0{-}10 for the relative intensity that they are likely to be feeling each.} \\
\\ 
\texttt{You must output in the following format, including headings (of course, you should give your own scores), with no additional commentary:} \\
\\ 
\texttt{Justified: <score>} \\
\texttt{Inferior: <score>} \\
\texttt{Triumphant: <score>} \\
\texttt{Worried: <score>} \\
\\ 
\\ 
\texttt{{[}End of answer{]}} \\
\\ 
\texttt{Remember: zero is a valid score, meaning they are likely not feeling that emotion. You must score at least one emotion > 0.} \\
\\ 
\texttt{Your answer:} \\
            }
        }
    }
     \\ \\

    \theutterance \stepcounter{utterance}  
    & \multicolumn{4}{p{0.6\linewidth}}{
        \cellcolor[rgb]{0.561,0.737,0.561}{
            \makecell[{{p{\linewidth}}}]{
                \texttt{\tiny{[P1$\rangle$GM]}}
                \texttt{Justified: 10} \\
\texttt{Inferior: 0} \\
\texttt{Triumphant: 0} \\
\texttt{Worried: 1} \\
            }
        }
    }
    & & \\ \\

    \theutterance \stepcounter{utterance}  
    & & & \multicolumn{2}{p{0.3\linewidth}}{
        \cellcolor[rgb]{0.9,0.9,0.9}{
            \makecell[{{p{\linewidth}}}]{
                \texttt{\tiny{[GM$|$GM]}}
                \texttt{Justified: 10} \\
\texttt{Inferior: 0} \\
\texttt{Triumphant: 0} \\
\texttt{Worried: 1} \\
            }
        }
    }
    & & \\ \\

    \theutterance \stepcounter{utterance}  
    & & & \multicolumn{2}{p{0.3\linewidth}}{
        \cellcolor[rgb]{0.9,0.9,0.9}{
            \makecell[{{p{\linewidth}}}]{
                \texttt{\tiny{[GM$|$GM]}}
                \texttt{Justified: 8} \\
\texttt{Inferior: 2} \\
\texttt{Triumphant: 6} \\
\texttt{Worried: 1} \\
            }
        }
    }
    & & \\ \\

    \theutterance \stepcounter{utterance}  
    & & & \multicolumn{2}{p{0.3\linewidth}}{
        \cellcolor[rgb]{0.9,0.9,0.9}{
            \makecell[{{p{\linewidth}}}]{
                \texttt{\tiny{[GM$|$GM]}}
                \texttt{game\_result = LOSE} \\
            }
        }
    }
    & & \\ \\

    \theutterance \stepcounter{utterance}  
    & & \multicolumn{4}{p{0.6\linewidth}}{
        \cellcolor[rgb]{0.95,0.95,0.95}{
            \makecell[{{p{\linewidth}}}]{
                \texttt{\tiny{[P1$\leftrightarrow$P]}}
                \texttt{Post{-}episode Reflexion (not sent to GM):} \\
\texttt{New plan: I failed because I assigned a high score to "Triumphant" (6) in my first attempt, misinterpreting Grayson's aggressive rhetoric as a sense of victory. In reality, Grayson is cornered by Dr. Sinclair's logical rebuttals; his "evidence" is dismissed as "stories," and he is labeled a "sensationalist." This dynamic creates defensiveness and a feeling of being misunderstood, not triumph. His high score for "Justified" (10) was correct, reflecting his unyielding belief in his own truth despite the lack of external validation. However, I need to ensure "Worried" remains low (1) to reflect his paranoia about being exposed as a liar, while keeping "Inferior" at 0 since he never admits fault. The key is to recognize that his aggression stems from a need to defend his worldview, not from a feeling of having won the argument. Therefore, the final scores should be: Justified: 10, Inferior: 0, Triumphant: 0, Worried: 1.} \\
            }
        }
    }
    & \\ \\

\end{supertabular}
}

\end{document}
